\section{Тестирование}

\subsection{Сборка проекта CMake}

Файл \texttt{CMakeLists.txt} настраивает сборку статической библиотеки \texttt{polynomial}, исполняемого файла \texttt{main} и тестового исполняемого файла \texttt{test\_polynomial}. Google Test подтягивается автоматически через \texttt{FetchContent}.

\lstinputlisting[style=cmakestyle]{../../CMakeLists.txt}

\subsection{Результаты тестирования}

Тесты размещены в \texttt{tests/test\_polynomial.cpp} (27 тестов). Примеры проверяемых сценариев: конструирование из массива, обрезка старших нулевых (\texttt{trim}), обработка пустого/нулевого многочлена, глубокое копирование, перемещение, присваивание, печать, вычисление в точке, вычисление для массива точек, сложение и умножение многочленов.

Сборка и запуск тестов (из корня \texttt{lab\_2}):

\begin{verbatim}
cmake -S . -B build
cmake --build build
ctest --test-dir build --output-on-failure
\end{verbatim}

Все тесты успешно проходят.
